% Ad Familiares 7.11 --- headnote
% ad-familiares-07-11, January 53 BC, Rome (perhaps Tusculum)

\textit{Cicero to C. Trebatius Testa, written probably from his
Tusculan villa in January 53 BC. Trebatius, a young jurist whom
Cicero had commended to Caesar, was serving on the Gallic staff
and grumbling, as ever, about the discomforts of the camp. The
opening joke turns on the political vacuum at Rome: with the
consular elections obstructed through most of 53, the city
lurched from one short \textit{interregnum} to the next, and a
civil lawyer in Rome had no courts to plead in --- so, Cicero
quips, even if Trebatius had not already left, he would be
leaving now.}

\textit{The middle section is the heart of the letter. Cicero is
relieved to find Trebatius joking again in his correspondence
(``these signs are better than the statues in my Tusculan
villa''), but probes for hard news: is Caesar actually advancing
him, or merely consulting him? If something is in train, stay;
if not, come home. The threat of comic exposure --- the mimographer
Laberius, and a second writer named Valerius, putting ``the British
jurisconsult'' on stage --- belongs to the running joke of the
Trebatius correspondence, in which Cicero pretends that his
friend's Gallic adventure is fit material for the mime stage.
The closing paragraph drops the banter for a moment of plain
counsel, then returns to warmth: whatever Trebatius wants he will
get, by his own merit and Cicero's effort.}
